% from ega0-1-fr.tex
\oldpage[0\textsubscript{I}]{68}

grand, d'où $\mathfrak J^p\subset\mathfrak m$ et
$\mathfrak J\subset\mathfrak m$ puisque $\mathfrak m$ est premier. Enfin,
$a)$ entraîne $b)$~: soit en effet $P$ l'adhérence de $\{0\}$ dans un
$A$-module $M$ de type fini, pour la topologie $\mathfrak J$-préadique~; en
vertu du th. de Krull (7.3.2), la topologie induite sur $P$ par la topologie
$\mathfrak J$-préadique de $M$ est la topologie $\mathfrak J$-préadique de
$P$, donc $\mathfrak JP=P$~; comme $P$ est de type fini, il résulte du lemme
de Nakayama que $P=0$ ($\mathfrak J$ étant contenu dans le radical de $A$).

On notera que les conditions de (7.3.5) sont remplies lorsque $A$ est un
\emph{anneau local noethérien} et $\mathfrak J\neq A$ un idéal quelconque de
$A$.

\begin{corollary}[7.3.6]
\label{0.7.3.6-fr}
Si $A$ est un anneau $\mathfrak J$-adique noethérien, tout $A$-module de type
fini est séparé et complet pour la topologie $\mathfrak J$-préadique.
\end{corollary}

Comme on a alors $\widehat A=A$, cela résulte aussitôt de (7.3.3).

On en conclut que la prop. (7.2.9) donne la description de \emph{tous} les
modules de type fini sur un anneau adique noethérien.

\begin{corollary}[7.3.7]
\label{0.7.3.7-fr}
Sous les hypothèses de (7.3.2), le noyau de l'application canonique
$M\to\widehat M=M\otimes_A\widehat A$ est l'ensemble des $x\in M$ annulés par
un élément de $1+\mathfrak J$.
\end{corollary}

En effet, pour que $x\in M$ appartienne à ce noyau, il faut et il suffit que
le séparé complété du sous-module $Ax$ se réduise à 0 (par (7.3.2)), autrement
dit que $x\in\mathfrak Jx$.

\subsection{Modules quasi-finis sur les anneaux locaux}
\label{subsection:0.7.4-fr}

\begin{definition}[7.4.1]
\label{0.7.4.1-fr}
Étant donné un anneau local $A$, d'idéal maximal $\mathfrak m$, on dit qu'un
$A$-module $M$ est quasi-fini (sur $A$) si $M/\mathfrak mM$ est de rang fini
sur le corps résiduel $k=A/\mathfrak m$.
\end{definition}

Lorsque $A$ est noethérien, le séparé complété $\widehat M$ de $M$ pour la
topologie $\mathfrak m$-préadique est alors un
$\widehat A$-\emph{module de type fini}~; en effet, comme
$\mathfrak m/\mathfrak m^2$ est alors un $A$-module de type fini, cela résulte
de (7.2.12) et de l'hypothèse sur $M/\mathfrak mM$.

En particulier, si on suppose de plus que $A$ est \emph{complet} et $M$
\emph{séparé} pour la topologie $\mathfrak m$-préadique (autrement dit,
$\bigcap_n\mathfrak m^nM=0$), $M$ lui-même est un
$A$-\emph{module de type fini}~: en effet, $\widehat M$ est alors un
$A$-module de type fini, et comme $M$ s'identifie à un sous-module de
$\widehat M$, $M$ est lui aussi de type fini (et d'ailleurs identique à son
complété en vertu de (7.3.6)).

\begin{proposition}[7.4.2]
\label{0.7.4.2-fr}
Soient $A$, $B$ deux anneaux locaux, $\mathfrak m$, $\mathfrak n$ leurs idéaux
maximaux, et supposons $B$ noethérien. Soient $\varphi:A\to B$ un
homomorphisme local, $M$ un $B$-module de type fini. Si $M$ est un $A$-module
quasi-fini, les topologies $\mathfrak m$-préadique et $\mathfrak n$-préadique
sur $M$ sont identiques, donc séparées.
\end{proposition}

Remarquons que par hypothèse $M/\mathfrak mM$ est de \emph{longueur finie} en
tant que $A$-module, donc aussi \emph{a fortiori} en tant que $B$-module. On en
conclut que $\mathfrak n$ est le \emph{seul idéal premier} de $B$ contenant
l'annulateur de $M/\mathfrak mM$~: en effet, on se ramène aussitôt, en vertu
de (I.7.4) et (I.7.2), au cas où $M/\mathfrak mM$ est \emph{simple}, donc
nécessairement isomorphe à $B/\mathfrak n$, et notre assertion est évidente
dans ce dernier cas. D'autre part, comme $M$ est un $B$-module de type fini,
les idéaux premiers qui contiennent l'annulateur de $M/\mathfrak mM$ sont ceux
qui contiennent $\mathfrak mB+\mathfrak b$, en désignant par $\mathfrak b$
l'annulateur du $B$-module $M$ (I.7.5). Comme $B$ est noethérien, on en conclut
([III], p.~127, cor.~4) que $\mathfrak mB+\mathfrak b$ est un idéal

